\documentclass[11pt]{article}
\usepackage{amsmath, amssymb}
\usepackage[margin=1in]{geometry}
\usepackage{booktabs}
\usepackage{siunitx}
\usepackage{hyperref}

\title{Sherlock Plan: Global Grid LiCl Device-Parameter Sweep}
\author{}
\date{}

\begin{document}
\maketitle

\section*{Reference and scope}

This document describes the plan for \texttt{solar\_lumped/scripts/grid\_param\_sweep.py}: a
full-factorial device-parameter sweep at every land point of a \SI{3}{\degree} lat/lon grid,
using the Wilson lumped SAWH model with a LiCl salt bed. One invocation of the script runs the
full parameter grid for a single site; the Sherlock job is a SLURM array with one task per site.

\section{Fixed parameters}\label{sec:fixed-above}

\begin{table}[h]
\centering
\small
\begin{tabular}{@{}lll@{}}
\toprule
\textbf{Parameter} & \textbf{Value} & \textbf{Notes} \\
\midrule
Salt & LiCl & only candidate; no feasibility fallback to other salts \\
Salt:polymer ratio & 4.0 & \\
Insulation gap & \SI{5}{\milli\meter} & \\
Vapor gap & \SI{40}{\milli\meter} & fixed, not swept -- see Section~\ref{sec:swept} \\
Tilt & \SI{35}{\degree} & \\
Weather year & 2024 & Open-Meteo historical-forecast API, cached locally \\
Grid spacing & \SI{3}{\degree} & land points only (Natural Earth 110m polygons) \\
Excluded region & Canada, Russia above \SI{60}{\degree}N & country polygons, not just a \\
 & & latitude line -- Alaska/Greenland/Scandinavia unaffected \\
Grid site count & 1405 & lat $\in[-54,72]$ minus the excluded region; computed at run time \\
 & & via \texttt{grid\_land\_points} \\
\bottomrule
\end{tabular}
\end{table}

\section{Swept parameters}\label{sec:swept}

Four axes, $5\times3\times3\times3 = 135$ combinations per site. Vapor gap was considered as a
fifth axis but dropped from the sweep (fixed at the \SI{40}{\milli\meter} baseline instead --
Section~\ref{sec:fixed-above}) to keep the grid smaller. Two of the remaining swept parameters
(absorber emissivity, glass transmission) were not previously exposed by any script -- they are
hardcoded constants (\texttt{table\_s3.EPS\_ABS}, \texttt{table\_s3.TAU\_GLASS}) unless overridden
via an explicit \texttt{DeviceThermalParams}, which is how \texttt{grid\_param\_sweep.py} wires
them in.

\begin{table}[h]
\centering
\small
\begin{tabular}{@{}llll@{}}
\toprule
\textbf{Parameter} & \textbf{CLI flag} & \textbf{Values} & \textbf{Baseline} \\
\midrule
Hydrogel thickness (mm) & \texttt{--hydrogel-thickness-mm} & 1.0, 3.25, 5.5, 7.75, 10.0 & 4.0 \\
Absorber emissivity $\varepsilon_{abs}$ & \texttt{--eps-abs} & 0.85, 0.90, 0.95 & 0.95 \\
Glass transmittance $\tau_{glass}$ & \texttt{--tau-glass} & 0.80, 0.85, 0.90 & 0.90 \\
Condenser fin area ratio & \texttt{--fin-area-ratio} & 3.0, 7.1, 12.0 & 7.1 \\
\bottomrule
\end{tabular}
\end{table}

Total problem size: $1405 \times 135 = 189{,}675$ independent (site, combo) evaluations.

\section{Time resolution: monthly, not single mean-day}\label{sec:resolution}

\subsection{Decision}
The default (\texttt{--resolution monthly}) simulates 12 Aitken-converged representative mean
days per site per combo (one per calendar month, day-weighted average), \emph{not} a single
annual mean day. \texttt{--resolution single} (12$\times$ cheaper) remains available for quick
previews, but is \emph{not} used for the production sweep.

\subsection{Why}
This sweep feeds a map that must be accurate everywhere -- not a search that only needs to find
the best few sites, where errors at non-competitive sites could be waved away. That changes the
trade-off. Validated by simulating the real day-by-day sequential year (365/366 days, gel state
carried forward) against the single mean-day and monthly-mean-day shortcuts, at three test
climates, after fixing \texttt{find\_cyclic\_state} to actually converge to the steady periodic
state (Aitken $\Delta^2$ extrapolation) rather than a fixed, possibly under-converged, warmup
count:

\begin{table}[h]
\centering
\small
\begin{tabular}{@{}lrrrr@{}}
\toprule
\textbf{Site} & \textbf{Annual (kg/m\textsuperscript{2})} & \textbf{1 mean-day} & \textbf{12 months} & \textbf{52 weeks} \\
\midrule
Patagonian coast (high seasonal variance) & 0.3809 & 0.3010 (-21.0\%) & 0.3287 (-13.7\%) & 0.3321 (-12.8\%) \\
Atacama desert (low seasonal variance) & 1.9052 & 1.9018 (-0.2\%) & 1.9011 (-0.2\%) & 1.9007 (-0.2\%) \\
Cambridge, MA (moderate seasonal variance) & 0.9032 & 0.7393 (-18.1\%) & 0.8606 (-4.7\%) & not run \\
\bottomrule
\end{tabular}
\end{table}

A single mean day is off by double digits ($-21\%$, $-18\%$) at exactly the sites -- high
seasonal variance -- where a map most needs to be trustworthy; monthly closes most of that gap
(to $-14\%$ and $-5\%$) for $12\times$ the compute, which is the right price for a map product.
Weekly buys almost nothing further over monthly (another $4.3\times$ cost for $\lesssim 1$
percentage point at Patagonia, confirmed twice); it is not used. Monthly is the sweet spot:
close to full-annual accuracy everywhere, at a small fraction of full-annual's cost
($\sim$33$\times$ cheaper -- Section~\ref{sec:cost}).

If a specific region still needs full-annual-level accuracy after seeing the monthly map (e.g. a
published headline number, or a site under closer scrutiny), rerun just that site with
\texttt{--resolution single} swapped for a direct call to
\texttt{scripts/compare\_annual\_vs\_mean\_day.py} (Approach A), which is affordable one site at a
time even though it is not affordable across the whole grid.

\subsection{Numerics: Aitken convergence and the 2-cycle fallback}
The steady periodic post-desorption state $(c_w, H)$ for a repeated day is found by
\texttt{find\_cyclic\_state} (restarted vector Aitken $\Delta^2$ extrapolation), typically
converging in 3--6 rounds (2 real one-day simulations per round) instead of the 100+ cycles
plain fixed-point warmup can need at strongly seasonal sites. Some (profile, config) pairs have
no single fixed point at all -- the one-cycle map bifurcates into a stable period-2 orbit (an
alternating wetter-day/drier-day pattern under the repeated forcing). This is detected as two
consecutive rounds where the relative step fails to shrink by at least half, and handled by
returning the average of the two alternating states (with a one-line log message) rather than an
arbitrary snapshot from whichever round happened to hit the round cap.

\section{What gets saved}

One CSV row per (site, combo), appended incrementally (not buffered), so a preempted job can
resume without redoing completed combos:

\begin{center}
\small
\texttt{lat, lon, mean\_rh\_frac, mean\_t\_amb\_c, mean\_solar\_w\_m2, salt,
hydrogel\_thickness\_mm, eps\_abs, tau\_glass, fin\_area\_ratio, vapor\_gap\_mm,
warmup\_method, resolution, mean\_yield\_kg\_m2, mean\_eta\_thermal, n\_periods}
\end{center}

\texttt{mean\_rh\_frac}, \texttt{mean\_t\_amb\_c}, and \texttt{mean\_solar\_w\_m2} are site
properties, not per-combo ones (identical across all 135 rows for a given site) -- day-weighted
means of the same monthly profiles already built for the ODE runs, so they cost nothing beyond
aggregating data already in memory. \texttt{mean\_solar\_w\_m2} is averaged over the desorption
(daylight) phase only; blending in the absorption (night) phase, which is by definition
low-solar, would dilute it with $\sim$12h of near-zero values.

Otherwise this is water yield and thermal efficiency only -- no LCOW or economics. That can be
computed from yield as a separate downstream step, the same way \texttt{lcow\_full\_global\_map.py}
does.

\section{Smoke test}\label{sec:smoke}

Validated end to end, both locally and on Sherlock itself, before any full-grid submission.

\begin{enumerate}
  \item \textbf{Plumbing check} (no ODE, instant): built device configs for two combos at the
  extremes of all four swept axes plus the fixed vapor gap and confirmed each reaches
  \texttt{DeviceConfig} / \texttt{DeviceThermalParams} as the intended value. Passed.
  \item \textbf{End-to-end run}, Atacama desert site (\ang{-23.6}, \ang{-70.4}), 2 combos
  (\texttt{eps\_abs} = 0.90 and 0.95, everything else at baseline), \texttt{--resolution monthly}
  (the production default):
\begin{verbatim}
python scripts/grid_param_sweep.py \
  --lat-lon -23.6 -70.4 \
  --hydrogel-thickness-mm 4.0 --eps-abs 0.90 0.95 \
  --tau-glass 0.85 --fin-area-ratio 7.1 \
  --output-csv outputs/smoke_test.csv
\end{verbatim}
  Result: combo 1 (eps\_abs=0.90) $\to$ yield \SI{1.707476}{\kilo\gram\per\meter\squared} in
  \SI{386.9}{\second}; combo 2 (eps\_abs=0.95) $\to$ yield
  \SI{1.800478}{\kilo\gram\per\meter\squared} in \SI{318.2}{\second}. Higher absorber emissivity
  correctly gives higher yield and higher thermal efficiency -- the direction physics predicts,
  confirming the full chain (weather fetch $\to$ 12 monthly mean-day profiles $\to$ Aitken steady
  state per month, day-weighted $\to$ yield) is wired correctly, not just that the code runs
  without error. The 2-cycle fallback (Section~\ref{sec:resolution}) fired on 3 of the 24
  month-evaluations across the two combos, each logged with its stall value and round number.

\end{enumerate}

\subsection{Sherlock-side smoke test}
Environment setup hit one snag worth documenting: the default \texttt{uv pip install} tried to
build \texttt{pyproj} from source (no system PROJ library on this node) because the resolved
\texttt{cartopy}/\texttt{pyproj} versions' wheels target a newer glibc than this cluster ships.
Forcing binary-only installs resolves to slightly older, wheel-compatible versions instead of
building anything:

\begin{verbatim}
ml python/3.12.1 uv
uv venv && source .venv/bin/activate
uv pip install --only-binary :all: numpy scipy pandas requests-cache retry-requests shapely cartopy

sbatch <<'EOF'
#!/bin/bash
#SBATCH --job-name=sawh-smoke
#SBATCH --time=01:00:00
#SBATCH --cpus-per-task=1
#SBATCH --mem-per-cpu=4G
#SBATCH --partition=serc
source .venv/bin/activate
export PYTHONUNBUFFERED=True
python scripts/grid_param_sweep.py \
  --lat-lon -23.6 -70.4 \
  --hydrogel-thickness-mm 4.0 --eps-abs 0.90 0.95 \
  --tau-glass 0.85 --fin-area-ratio 7.1 \
  --output-csv outputs/smoke_test_sherlock.csv --resume
EOF
\end{verbatim}

(\texttt{--time=01:00:00}, not \SI{20}{\minute} as first tried -- see below for why.
\texttt{--resume} makes a resubmission after a timeout/preemption pick up where it left off
instead of redoing the first combo.)

\textbf{Result: yields matched the local run exactly} --
\SI{1.707476}{\kilo\gram\per\meter\squared} and \SI{1.800478}{\kilo\gram\per\meter\squared},
confirming the full chain (weather fetch $\to$ 12 monthly mean-day profiles $\to$ Aitken steady
state per month, day-weighted $\to$ yield) runs identically on the cluster, not just that it
runs without error.

\textbf{But it ran much slower than the local Mac}: combo 1 took \SI{1011.1}{\second} on
Sherlock versus \SI{386.9}{\second} locally; combo 2 took \SI{859.9}{\second} versus
\SI{318.2}{\second} locally ($2.61$--$2.70\times$, average $2.66\times$). The first submission
used \texttt{--time=00:20:00} (sized off the Mac timing) and got killed by the wall-clock limit
mid-combo-2; resubmitting with \texttt{--resume} and more time completed it. This
$2.66\times$ factor is what Section~\ref{sec:cost}'s cost estimate is actually calibrated on --
a laptop benchmark alone would have understated the real job by the same factor.

\section{Full sweep}

\subsection{Job structure}
\textbf{One array task per site running all 135 combos serially does not work}: at the measured
per-combo cost (Section~\ref{sec:cost}), one site takes 37.9--131.9 \emph{hours} serially on
Sherlock, far beyond any reasonable single-task \texttt{--time}. Each site's 135 combos are
instead split into 45 chunks of 3 combos (135 divides evenly), so each array task is one (site,
chunk) pair costing roughly \SI{0.84}{}--\SI{2.93}{\hour} on Sherlock (see
Section~\ref{sec:cost} for how that's calibrated). \texttt{--combo-offset}/\texttt{--combo-limit}
(added to \texttt{grid\_param\_sweep.py} for this) slice the site's combo grid; the two
array-index components are recovered from a single flattened SLURM index via integer div/mod:

\begin{verbatim}
#!/bin/bash
#SBATCH --job-name=sawh-sweep
#SBATCH --array=0-63224%300
#SBATCH --time=05:00:00
#SBATCH --cpus-per-task=1
#SBATCH --mem-per-cpu=4G
#SBATCH --partition=owners

N_CHUNKS_PER_SITE=45
CHUNK_SIZE=3
SITE_INDEX=$(( SLURM_ARRAY_TASK_ID / N_CHUNKS_PER_SITE ))
CHUNK_INDEX=$(( SLURM_ARRAY_TASK_ID % N_CHUNKS_PER_SITE ))
COMBO_OFFSET=$(( CHUNK_INDEX * CHUNK_SIZE ))

source .venv/bin/activate
export PYTHONUNBUFFERED=True
python scripts/grid_param_sweep.py \
  --site-index $SITE_INDEX --step 3 \
  --combo-offset $COMBO_OFFSET --combo-limit $CHUNK_SIZE \
  --output-csv outputs/grid_sweep/site_${SITE_INDEX}_chunk_${CHUNK_INDEX}.csv --resume
EOF
\end{verbatim}

\texttt{\%300} throttles to 300 concurrent tasks regardless of how many are queued; raise or
lower it based on how the partition responds. Each (site, chunk) gets its own output file --
\emph{not} one shared file per site -- because multiple chunks of the same site can run
concurrently on different nodes, and concurrent appends/header-writes to one shared file would
race. \texttt{--resume} makes a preempted/resubmitted task skip combos already written in its
own chunk file, so preemption on \texttt{owners} costs at most the in-progress combo, not the
whole chunk. \texttt{--time=05:00:00} leaves nearly $1.7\times$ margin over the slowest measured
chunk (\SI{2.93}{\hour}); the smaller 3-combo chunk (versus the originally-planned 9) trades more
array tasks for a shorter preemption-exposure window per task, now that Sherlock's real per-combo
cost is known to run higher than the local estimate.

\texttt{MaxArraySize} on this cluster was checked directly (\texttt{scontrol show config | grep
-i MaxArraySize}) and is \textbf{1{,}000{,}000} -- 63,225 array elements fits with enormous
headroom, so no splitting across multiple \texttt{sbatch} calls is needed.

\subsection{Weather prefetch}\label{sec:weather}
Each site needs one Open-Meteo fetch (cached permanently now that \texttt{WeatherClient}'s cache
TTL is fixed -- see repository history). 1405 concurrent Sherlock array tasks must \emph{not}
each independently hit the API: prefetch/cache all 1405 sites in one throttled pass first.

Originally planned as a local (Mac) prefetch followed by an \texttt{rsync} of the resulting
21GB \texttt{.weather\_cache} to \texttt{\$GROUP\_HOME}. In practice, home-network upload
bandwidth ($\sim$\SI{1.3}{\mega\byte\per\second} measured) made that transfer impractically
slow (multi-hour estimate, and it dropped mid-transfer at least once). \textbf{Run the prefetch
directly on Sherlock instead} -- there's no reason to build the cache locally and ship it when
Sherlock's own network link to Open-Meteo is unrelated to a laptop's home upload speed:

\begin{verbatim}
sh_dev
cd $GROUP_HOME/SAWH_TEAs/solar_lumped
source .venv/bin/activate
python scripts/prefetch_weather.py --step 3
\end{verbatim}

Use an interactive \texttt{sh\_dev} session rather than the login node (light on CPU but runs
for several minutes, which is what \texttt{sh\_dev} is for). \texttt{prefetch\_weather.py} builds
the land/country-exclusion geometry and imports the physics stack once, then fetches every site
in the same process and connection pool ($\sim$6$\times$ faster than a subprocess-per-site
loop), and retries failed fetches with an exponential cooldown (Open-Meteo's free tier is
rate-limited to 600 calls/min, 5{,}000/hour, 10{,}000/day) instead of giving up on the first 429.
With this cache now living directly on Sherlock, drop \texttt{--cache-dir} from the sweep
commands below -- the script's default (relative to the repo) already resolves correctly.

\subsection{Merging outputs}\label{sec:merge}
The array produces up to 63,225 separate CSV files (one per (site, chunk), all sharing the same
header). Concatenate them into one dataset for mapping once the array finishes:

\begin{verbatim}
n=$(ls outputs/grid_sweep/site_*_chunk_*.csv 2>/dev/null | wc -l)
echo "found $n of 63225 expected chunk files"

{ head -1 "$(ls outputs/grid_sweep/site_*_chunk_*.csv | head -1)"
  for f in outputs/grid_sweep/site_*_chunk_*.csv; do tail -n +2 "$f"; done
} > outputs/grid_sweep_full.csv
\end{verbatim}

Check the file count before treating the merge as final -- a short count means some array tasks
haven't completed (still queued, still running, or failed and need resubmitting), not that the
sweep silently skipped anything. Row count in the merged file should be at most $189{,}675$
(fewer only if some chunks are still outstanding; \texttt{--resume} prevents it from ever being
more).

\subsection{Expected cost}\label{sec:cost}
Mac-measured per-combo cost at \texttt{--resolution monthly} (post 2-cycle fix): \SI{380.0}{\second}
(Atacama desert, low-seasonal-variance) to \SI{1323.4}{\second} (Cambridge, moderate-seasonal-
variance). \textbf{Sherlock itself runs $\sim$2.66$\times$ slower per combo than that Mac} for
this workload -- measured directly from the Sherlock smoke test (Section~\ref{sec:smoke}):
\SI{1011.1}{\second} and \SI{859.9}{\second} there versus \SI{386.9}{\second} and
\SI{318.2}{\second} locally for the identical two combos (ratios 2.61$\times$ and 2.70$\times$).
Scaling the Mac-measured fast/slow bounds by that factor gives the real Sherlock-calibrated cost.
Total: $189{,}675$ units (1405 sites $\times$ 135 combos).

\begin{table}[h]
\centering
\small
\begin{tabular}{@{}lrrrr@{}}
\toprule
\textbf{Cost basis} & \textbf{Core-hours} & \textbf{@500 cores} & \textbf{@1000 cores} & \textbf{@2000 cores} \\
\midrule
Fast/typical (\SI{1010}{\second}/combo) & 53{,}214 & 106.4 h & 53.2 h & 26.6 h \\
Slow (\SI{3518}{\second}/combo) & 185{,}324 & 370.6 h & 185.3 h & 92.7 h \\
\bottomrule
\end{tabular}
\end{table}

So: realistically \textbf{53{,}000--185{,}000 core-hours}, i.e. roughly 1--8 days at 1000
concurrent cores. This is $\sim$2.66$\times$ the earlier Mac-only estimate -- the lesson being
that a laptop benchmark is a starting point, not a substitute for measuring on the actual
target hardware before sizing the real job. Dropping vapor gap from the sweep
(Section~\ref{sec:swept}, $405\to135$ combos) and excluding northern Canada/Russia
(Section~\ref{sec:fixed-above}, $1635\to1405$ sites) still meaningfully reduce this from what it
would otherwise have been; monthly over single (Section~\ref{sec:resolution}) costs $12\times$
more than the cheapest option, but buys the accuracy the map actually needs.

\subsection{No GPU}
This code is scalar SciPy (\texttt{solve\_ivp}, Radau/BDF methods) -- single-threaded CPU work,
no GPU path exists. Target CPU-only allocation on \texttt{serc}/\texttt{owners}; do not request
A100 nodes for this job.

\section{Open items / caveats}
\begin{itemize}
  \item Per-site cost varies by climate; the $3.5\times$ spread above is from only 2 sampled
  sites out of 1405 -- treat the cost table as order-of-magnitude, not a tight bound.
  \item \texttt{find\_cyclic\_state}'s 2-cycle fallback logs a one-line warning whenever it
  fires (it fired 2--4 times per combo at the two sites tested here); the full sweep's logs give
  a free count of how often the grid actually hits this bifurcation, which was previously
  unknown.
  \item Vapor gap is fixed at the \SI{40}{\milli\meter} baseline rather than swept. If it turns
  out to matter for the map, add it back as a fifth axis at the cost of another $\sim3\times$
  (135$\to$405 combos).
  \item Northern Canada and Russia above \SI{60}{\degree}N are excluded by country polygon, not
  just a latitude line (Section~\ref{sec:fixed-above}) -- Alaska, Greenland, Iceland, and
  Scandinavia's far north are still included. Pass
  \texttt{exclude\_country\_above\_lat=\{\}} to \texttt{grid\_land\_points} to disable this and
  restore the full 1635-site grid if that exclusion turns out to be unwanted.
  \item If a specific site needs full-annual-level accuracy after seeing the monthly map, rerun
  it directly with \texttt{compare\_annual\_vs\_mean\_day.py} rather than resubmitting the whole
  grid at annual resolution ($\sim$33$\times$ monthly's cost, weeks of wall-clock even at large
  core counts).
\end{itemize}

\end{document}
